% 
%       قالب زیر توسط سید صالح اعتمادی جهت جزوه برداری توسط دانشجویان آماده شده است.
% 		استفاده از این قالب برای عموم با حفظ این دو خط در ابتدای فایل بلا مانع است.
%
% تشکر ویژه از آقای علی حیدری که در تهیه این قالب از قالب‌های لاتک تهیه شده توسط ایشان کمک گرفته شده

\documentclass[oneside]{book}

\def \CourseName {مقدمه ای بر برنامه نویسی بازی های سه بعدی با DirectX 12}
\def \Instructor {دکتر رضا روحانی}
\def \Semester {نیم‌سال اول\\سال تحصیلی 1402-1403}

\usepackage[]{algorithm2e}
\usepackage{cite}
\usepackage{calc}
\usepackage{fancyhdr}
\usepackage{lipsum}
\usepackage{color}
\usepackage{ragged2e}
\usepackage[inline]{enumitem}
\usepackage[dvipsnames]{xcolor}
\usepackage{graphicx}
\usepackage{wrapfig}
\usepackage{float}


\usepackage[skip=12pt,indent=2em]{parskip}
\usepackage{setspace}
\usepackage{textcomp}
\usepackage{etoolbox}
\usepackage{xpatch}


\usepackage{tabu}
\usepackage{hyperref}
\usepackage{minted}
\usepackage{xepersian} 
\settextfont[Path=./fonts/,
BoldFont={XB Yas Bd}, 
ItalicFont={XB Yas It},
BoldItalicFont={XB Yas BdIt},
Extension = .ttf]{XB Yas}

%---------------------------------define color-------------------------------
\definecolor{gray}{rgb}{0.4,0.4,0.4}
\definecolor{dkgreen}{rgb}{0,0.6,0}
\definecolor{gray}{rgb}{0.4,0.4,0.4}
\definecolor{mauve}{rgb}{0.58,0,0.82}
\definecolor{mygray}{rgb}{0.9,0.9,0.9}
\definecolor{myCyan}{rgb}{0,1,1}
\definecolor{navyBlue}{HTML}{0a065d}
\definecolor{aliceblue}{rgb}{0.94, 0.97, 1.0}
%----------------------------------------------------------------------------


\setminted{
    frame=lines,
    framesep=2mm,
    baselinestretch=1.0,
    xrightmargin=0in,
    xleftmargin=0in,
    fontsize=\footnotesize,
    bgcolor=aliceblue,
    linenos=true
}

\hypersetup{
    colorlinks=true,
    linkcolor=blue,
    filecolor=magenta,      
    urlcolor=blue,
}
\setlist[enumerate]{itemsep=-3mm}
\setlist[itemize]{itemsep=-3mm}


\newfloat{program}{thp}{lop}
\floatname{program}{نمونه کد}
\newcommand{\programref}[1]{نمونه کد \ref{#1}}

\renewcommand{\chaptername}{جلسه}
\renewcommand{\baselinestretch}{1.5}
\newcommand{\grayBox}[1]{\colorbox{gray!10}{\lr{\texttt{#1}}}}
\newcommand{\grayLBox}[1]{\fbox{\grayBox{\parbox{11.8 cm}{#1}}}}

\AtBeginEnvironment{minted}{\setlength{\parskip}{-20pt}}
\xpretocmd{\inputminted}{\par\vspace{-10pt}}{}{}
\xapptocmd{\inputminted}{\par\vspace{-20pt}}{}{}


\makeatletter
\newcommand{\chapterauthor}[1]{%
    {\parindent0pt\vspace*{-25pt}%
        \linespread{1.1}\large\scshape#1%
        \par\nobreak\vspace*{35pt}}
    \@afterheading%
}
\makeatother
\title{
    \center
    \includegraphics[width=5cm, height=5.8cm]{images/SKU_logo_color.jpg} \\
    دانشکده فنی مهندسی \\[25pt]     
جزوه درس \\
\CourseName
}

\author{
    استاد درس:
    \Instructor \\[25pt]
}
\date{\Semester}

\begin{document}
    \pagenumbering{gobble}
    \pagestyle{headings}
    \maketitle

    %--------------------------------------%
\newpage

\textbf{\vspace{80pt}}
\title{
    \BTitr
    \center \Huge
    \begin{flushright}
        \textbf{قبل از شروع ...}
    \end{flushright}
}
\textbf{\vspace{80pt}}

{
    \Large
    کتابی پیش رو دارید، مباحث کتاب \lr{Introduction to 3d Game Programming With Directx 12} را در برمیگرد.
    در این کتاب از ترجمه‌ی محض خارج شده و موارد را به صورت گسترده‌تر بررسی می‌کند و هر مرحله را قدم به قدم توضیح می‌دهد.
    در اوایل ترجمه‌ی کتاب، نظر به این بود تا کتاب ترجمه شود ولی با توجه به اینکه منابع فارسی برای مباحث کتاب بسیار کم بوده و حتی نایاب است، تصمیم بر آن شد که اطلاعات و مباحث اضافی نیز به آن اضافه گردد.
    همچنین در تمامی کتاب سعی شده تا موارد تا جای ممکن به اصل خود تشابه داشته باشند.
    در هر صورت تمامی منابع استفاده شده و همه موارد مورد نیاز این منابع به صورت یکجا برای شما در لینک زیر قرار داده شده است.
    می‌توانید با اسکن \lr{QR Code} زیر نیز، آن‌ها را مشاهده کنید.

    \begin{figure}[H]
        \centering
        \setlength{\belowcaptionskip}{-10pt}
        \includegraphics[scale=0.15]{Images/2.BeforeIntro.1.1.png}
        \caption*{\Large \lr{\href{https://alibadiee.ir/Book/3d-DirectX12}{https://alibadiee.ir/Book/3d-DirectX12}}}
    \end{figure}

    این کتاب به صورت رایگان منتشر می‌شود و بجز برای خرید نسخه‌ی چاپی آن، شما نیاز به پرداخت هزینه‌ای ندارید.
    همچنین در صورتی که تمایل دارید، می‌توانید در بهتر کردن این اثر به ما کمک کنید. برای دریافت راهنما از لینک یا \lr{QR Code} زیر استفاده کنید.

    \begin{figure}[H]
        \centering
        \Large
        \setlength{\belowcaptionskip}{-10pt}
        \includegraphics[scale=0.15]{Images/2.BeforeIntro.1.2.png}
        \caption*{\Large \lr{\href{https://alibadiee.ir/Book/3d-DirectX12/Help-Editing}{https://alibadiee.ir/Book/3d-DirectX12/Help-Editing}}}
    \end{figure}
}

%--------------------------------------%

    \title{
        \center \Huge
        معرفی \\[25pt]
    }

    {Direct3D 12 یک کتابخانه رندر برای نوشتن برنامه های گرافیکی سه بعدی با کارایی بالا با استفاده از سخت افزار گرافیکی مدرن بر روی پلتفرم های مختلف ویندوز 10 (Windows Desktop، Mobile و Xbox One) است. Direct3D یک کتابخانه سطح پایین است به این معنا که رابط برنامه نویسی کاربردی آن (API) سخت افزار گرافیکی زیرینی را که کنترل می کند مدل سازی می کند. مصرف کننده اصلی Direct3D صنعت بازی است که در آن موتورهای رندر سطح بالاتر بر روی Direct3D ساخته می شوند. با این حال، صنایع دیگر به گرافیک سه بعدی تعاملی با کارایی بالا نیز نیاز دارند، مانند تجسم پزشکی و علمی و بررسی های معماری. علاوه بر این، با مجهز شدن هر رایانه شخصی جدید به یک کارت گرافیک مدرن، برنامه های غیر سه بعدی شروع به استفاده از GPU (واحد پردازش گرافیکی) برای تخلیه کار به کارت گرافیک برای محاسبات فشرده می کنند. این به عنوان محاسبات GPU با هدف عمومی شناخته می شود و Direct3D API سایه زن محاسباتی را برای نوشتن برنامه های GPU با هدف عمومی ارائه می دهد. اگرچه Direct3D 12 معمولاً از C++ بومی برنامه‌ریزی می‌شود، تیم SharpDX (http://sharpdx.org/) در حال کار بر روی Wrapper‌های .NET هستند تا بتوانید از برنامه‌های مدیریت شده به این API گرافیکی سه بعدی قدرتمند دسترسی داشته باشید.}

    {این کتاب مقدمه ای بر برنامه نویسی گرافیک های کامپیوتری تعاملی با تاکید بر توسعه بازی با استفاده از Direct3D 12 ارائه می دهد. اصول برنامه نویسی Direct3D و shader را آموزش می دهد و پس از آن خواننده آماده می شود تا تکنیک های پیشرفته تری را یاد بگیرد. کتاب به سه بخش اصلی تقسیم شده است. بخش اول ابزارهای ریاضی را توضیح می دهد که در سراسر این کتاب استفاده خواهد شد. بخش دوم نحوه پیاده سازی وظایف اساسی در Direct3D مانند مقداردهی اولیه را نشان می دهد. تعریف هندسه سه بعدی؛ راه اندازی دوربین ها؛ ایجاد رئوس، پیکسل، هندسه، و محاسبه سایه زن. نورپردازی؛ بافت سازی؛ مخلوط کردن؛ شابلون سازی; و تسلیت. بخش سوم عمدتاً در مورد استفاده از Direct3D برای پیاده سازی انواع تکنیک های جالب و جلوه های ویژه است، مانند کار با مش شخصیت های متحرک، برداشتن، نگاشت محیط، نقشه برداری معمولی، سایه های بلادرنگ، و انسداد محیط.}

    {برای مبتدیان، این کتاب بهتر است جلو به عقب بخواند. فصل ها به گونه ای سازماندهی شده اند که با هر فصل، مشکل به تدریج افزایش می یابد. به این ترتیب، هیچ جهش ناگهانی در پیچیدگی وجود ندارد که خواننده را گم کند. به طور کلی، برای یک فصل خاص، از تکنیک ها و مفاهیمی که قبلاً توسعه داده شده است استفاده خواهیم کرد. بنابراین، مهم است که قبل از ادامه، بر مطالب یک فصل تسلط داشته باشید. خوانندگان با تجربه می توانند فصل های مورد علاقه خود را انتخاب کنند. در نهایت، ممکن است از خود بپرسید که پس از خواندن این کتاب چه نوع بازی هایی را می توانید توسعه دهید. پاسخ این سوال را بهتر است با مرور این کتاب و مشاهده انواع برنامه های توسعه یافته به دست آورید. از این رو باید بتوانید انواع بازی هایی را که می توان بر اساس تکنیک های آموزش داده شده در این کتاب و برخی از نبوغ خود توسعه داد، تجسم کنید.}

    %--------------------------------------%
    \title{
        \LARGE
        مخاطب مورد نظر
    }
    \\ \rule{\textwidth}{0.5pt}
    {این کتاب با در نظر گرفتن سه مخاطب زیر طراحی شده است:}
    \begin{enumerate}
        \item {برنامه نویسان سطح متوسط C++ که می خواهند مقدمه ای بر برنامه نویسی سه بعدی با استفاده از آخرین نسخه}
        \item {برنامه نویسان سه بعدی با یک API غیر از DirectX (به عنوان مثال، OpenGL) تجربه کرده اند و می خواهند یک مقدمه برای Direct3D 12 داشته باشند.}
        \item {برنامه نویسان با تجربه Direct3D که مایل به یادگیری آخرین نسخه Direct3D هستند}
    \end{enumerate}

    \title{
        \LARGE
        پیش نیازها
    }
    {لازم به ذکر است که این مقدمه ای بر Direct3D 12، برنامه نویسی سایه زن و برنامه نویسی بازی های سه بعدی است. این مقدمه ای برای برنامه نویسی کامپیوتر عمومی نیست. خواننده باید شرایط زیر را رعایت کند:}

    \begin{enumerate}
        \item {ریاضیات دبیرستان: مثلاً جبر، مثلثات و توابع (ریاضی).}
        \item {صلاحیت با ویژوال استودیو: باید بداند که چگونه پروژه ها را ایجاد کند، فایل ها را اضافه کند و کتابخانه های خارجی را برای پیوند مشخص کند.}
        \item {C++ متوسط ​​و مهارت های ساختار داده: به عنوان مثال، با اشاره گرها، آرایه ها، بارگذاری بیش از حد اپراتورها، لیست های پیوندی، وراثت و چندشکلی راحت است.}
        \item {آشنایی با برنامه نویسی ویندوز با Win32 API مفید است، اما لازم نیست. ما یک پرایمر Win32 را در ضمیمه A ارائه می دهیم.}
    \end{enumerate}

    \title{
        \LARGE
        پیش نیازها
    }
    {}
    \begin{enumerate}
        \item {ریاضیات دبیرستان: مثلاً جبر، مثلثات و توابع (ریاضی).}
        \item {صلاحیت با ویژوال استودیو: باید بداند که چگونه پروژه ها را ایجاد کند، فایل ها را اضافه کند و کتابخانه های خارجی را برای پیوند مشخص کند.}
        \item {C++ متوسط ​​و مهارت های ساختار داده: به عنوان مثال، با اشاره گرها، آرایه ها، بارگذاری بیش از حد اپراتورها، لیست های پیوندی، وراثت و چندشکلی راحت است.}
        \item {آشنایی با برنامه نویسی ویندوز با Win32 API مفید است، اما لازم نیست. ما یک پرایمر Win32 را در ضمیمه A ارائه می دهیم.}
    \end{enumerate}

    \newpage
    %--------------------------------------%

    \pagenumbering{arabic}

    
\def \Subject {جبر برداری}
\def \Author {علی بدیعی}
\def \Date {1402/4/22}
\def \Session {1}
\setcounter{chapter}{\Session-1}

\part{پیشنیاز های ریاضی} \chapter{\Subject}
\chapterauthor{\Author~ - \Date}

%--------------------------------------%
\rule{\textwidth}{0.5pt}

\def \Subject {جبر ماتریسی}
\def \Author {علی بدیعی}
\def \Date {1402/4/22}
\def \Session {2}
\setcounter{chapter}{\Session-1}

\chapter{\Subject}
\chapterauthor{\Author~ - \Date}

%--------------------------------------%
\rule{\textwidth}{0.5pt}

\def \Subject {انتقال ها}
\def \Author {علی بدیعی}
\def \Date {1402/4/22}
\def \Session {3}
\setcounter{chapter}{\Session-1}

\chapter{\Subject}
\chapterauthor{\Author~ - \Date}

%--------------------------------------%
\rule{\textwidth}{0.5pt}

جزوه جلسه
\Session ام
 مورخ 
\Date
درس 
\CourseName
 تهیه شده توسط
 \Author. 


 
\section{معیار ارزیابی جزوه}
معیارهای مورد استفاده برای ارزشیابی کیفیت جزوه به شرح زیر است:
\begin{itemize}
    \item    
     پوشش کامل مطالب
    \item    
     رعایت قواعد نگارشی دستور زبان فارسی
    \item    
     استفاده از اشکال مناسب
    \item    
     اشاره به منابع کمک آموزشی
\end{itemize}


برای نوشتن شبه‌کد
 \footnote{pseudocode}
 می‌توانید از مثال زیر در الگوریتم
  \ref{pseudocode1}
   استفاده کنید :

\begin{latin}
\begin{algorithm}[H]
    \KwData{this text}
    \KwResult{how to write algorithm with \LaTeX2e }
    initialization\;
    \While{not at end of this document}{
        read current\;
        \eIf{understand}{
            go to next section\;
            current section becomes this one\;
        }{
            go back to the beginning of current section\;
        }
    }
    \label{pseudocode1}
    \caption{How to write algorithms}
\end{algorithm}
\end{latin}

همچنین برای اضافه کردن شکل می‌توانید از شکل زیر استفاده کنید و برای ارجاع دادن به بصورت شکل
 \ref{fig:dynamicprogramming}
 استفاده کنید.
 همچنین برای درج کلمات انگلیسی در پاراگراف فارسی می‌توان به این شکل
 \lr{English Text}
 یا برای تاکید به این شکل
 \grayBox{English Text}
 عمل کرد.
\begin{figure}[H]
    \centering
    \includegraphics[width=0.4\linewidth]{images/SKU_logo_color.jpg}
    \caption{برنامه نویسی پویا}
    \label{fig:dynamicprogramming}
\end{figure}

برای گنجاندن قطعه‌ای از کد به زبان‌های 
\grayBox{Java, C\#, C++, Python} 
از فایل اصلی کد می‌توانید به شکل زیر عمل کنید.
 
\begin{program}[H]
    \inputminted[frame=none]{csharp}{Sample.cs}
    \caption{تابع فیبوناچی در سی‌شارپ}
    \label{csharpfib}
\end{program}

به عنوان مثال در 
\programref{csharpfib}
پیاده‌سازی تابع فیبوناچی به زبان سی‌شارپ را ملاحظه می‌کنید. این تابع ورودی $n$ را دریافت کرده و عدد $n$ام دنباله را برمی‌گرداند. همانطور که ملاحظه می‌کنید این تابع به صورت بازگشتی پیاده‌سازی شده است.

\begin{program}[H]
\inputminted[frame=none]{cpp}{Sample.cpp}
\caption{تابع فیبوناچی در \grayBox{C++} }
\label{cppfib}
\end{program}

با مقایسه متن لاتک و خروجی به این مساله پی‌ خواهید برد که برای اشاره به یک قطعه کد در لاتک لازم است از دستور
\grayBox{label}
داخل محیط/بخش 
\grayBox{program}
استفاده کرده و برای اشاره به آن بخش از کد از دستور
\grayBox{programref}
استفاده نمایید. لازم است پارامتر داده شده به هر دو دستور یکی باشد تا بدرستی به قطعه کد مورد نظر اشاره کنید(\programref{pythonfib}). 

\begin{program}[H]
\inputminted[frame=none]{python}{sample.py}
\caption{تابع فیبوناچی در پایتون }
\label{pythonfib}
\end{program}

لطفا در همه موارد به جزئیات رعایت شده در متن لاتک دقت کنید. مخصوصا برای راحتی ویرایش متن لاتک بهتر است دستورات لاتین لاتک در خطوط جداگانه از متن فارسی قرار داده شود (\programref{javafib}).

\begin{program}[H]
\inputminted[frame=none]{java}{Sample.java}
\caption{تابع فیبوناچی در جاوا }
\label{javafib}
\end{program}
 چنانچه بخواهید قطعه کد را مستقیم در متن جای دهید، می‌توانید از الگوی زیر پیروی کنید (\programref{csharpnode}).
 
\begin{program}[H]
\begin{minted}[frame=none]{csharp}
public class Node<_Type>
{
    public _Type Value;
    public Node<_Type> Next;
}
\end{minted}
\caption{تعریف لیست پیوندی در سی‌شارپ}
\label{csharpnode}
\end{program}
 
همانطور که در نمونه‌کدها ملاحظه می‌کنید فاصله‌ها دقیقا همانطور که در لاتک آمده است در خروجی نمایش داده می‌شود.
 
\begin{program}[H]
\begin{minted}[frame=none]{java}
public class Node<_Type> {
    public _Type Value;
    public Node<_Type> Next;
}
\end{minted}
\label{javanode}
\caption{تعریف لیست پیوندی در ‌جاوا}
\end{program}

لذا لازم است که فاصله‌های ابتدای خط را مرتب و منظم و بدون کوچکترین اشکالی رعایت کنید تا خروجی مستند شما نیز بدون اشکال باشد.

\begin{program}[H]
\begin{minted}[frame=none]{cpp}
template<typename _Type>
class Node
{
public:
    _Type Value;
    Node<_Type>* Next;
}
\end{minted}
\label{cppnode}
\caption{تعریف لیست پیوندی در ‌\grayBox{C++}}
\end{program}
در نهایت استفاده از دستور 
\grayBox{grayBox}
نیز می‌تواند کمک شایانی به زیبایی و خوانایی متن شما بکند. این دستور علاوه بر عوض کردن رنگ پس‌زمینه از فونت انگلیسی با عرض ثابت نیز استفاده می‌کند که برای کلمات کلیدی یا اسامی متغیرها در کد قابل استفاده است.

\begin{program}[H]
\begin{minted}[frame=none]{python}
class Node: 
def __init__(self, data): 
    self.value = data
    self.next = None  
\end{minted}
\label{pythonnode}
\caption{تعریف لیست پیوندی در ‌پایتون}
\end{program}

\grayLBox{
    AAC\\
    ACG\\
    GAA\\
    GTT\\
    TCG\\
}
\\
همچنین سعی کنید حتی‌الامکان به منابع و مراحع مناسب ارجاع دهید
\cite{CLRS}. 
علاوه بر مراجع چنانچه ابزار یا وب‌سایت قابل توجهی موجود است خوب است به آن هم ارجاع دهید
\cite{visualizationwebsite}. 



    \begin{latin}
        \bibliography{references}
        \bibliographystyle{ieeetr}
    \end{latin}

\end{document}
